\section{Introduction}
\label{sec:intro}

DeepSeek-V4-Flash and V4.1-Flash store every routed expert in MXFP4~\citep{rouhani2023mx,deepseekv4}, which encodes
the weights as 4-bit E2M1 codes $\{0,\pm0.5,\pm1,\pm1.5,\pm2,\pm3,\pm4,\pm6\}$ with one power-of-two scale per 32
inputs, or 4.25 bits per weight. Because the experts hold almost all of the parameters, serving these models in less
memory requires going below the native format. For most of these models no higher-precision checkpoint is published,
and since the experts are trained with MXFP4 in the loop, such a checkpoint may not exist at
all~\citep{deepseekv4,kimik3,mimov25fp4}. Current practice nonetheless treats the MXFP4 weights as continuous.
EXL3~\citep{exllamav3}, a streamlined implementation of QTIP~\citep{tseng2024qtip}, dequantizes the experts, rotates
them with random Hadamard transforms until they look i.i.d.\ Gaussian, and codes them with a Gaussian-matched trellis.
GGUF builds below 4 bits instead re-round the dequantized weights onto llama.cpp's block grids.

Rotation is an appropriate prior for BF16 weights but not for MXFP4 weights. Every nonzero MXFP4 weight in a column
is an E2M1 value times a power of two, and because E2M1 is nearly closed under doubling and halving, the weights of a
column lie on a small union of grids. Such a source has finite entropy (\final{ENTROPY-RANGE} bits per weight on
DeepSeek-V4.1-Flash), and its memoryless distortion--rate function lies far below that of a Gaussian source, the
curve that Gaussian-matched coders are designed for. Rotation preserves the grid structure but moves it into the joint
distribution of each 128-coordinate Hadamard block, where a coder that models the rotated coordinates as i.i.d.\
Gaussian does not exploit it.

GRAFT keeps the parts of EXL3 that do not depend on rotation (the 16-bit bitshift trellis, the tile layout, the
hashed state-to-value map and the error-feedback encoder) but decodes to MXFP4 blocks. It drops the rotation
and maps EXL3's hash through one fused multiply-add and the GPU's E2M1 conversion, so that every 32-input block of a
column decodes to E2M1 values times one power of two. Each block carries a one-bit flag, which is the block's MX scale
and halves the block where the source sits an octave lower. A small band of states at the hash centre decodes to a
power-of-two ladder for outlier weights, and the decoder applies these values as an exact sparse correction. Columns
are scaled by a power of two and a fitted fp16 factor, and every matrix is coded at the same rate of 3 bits per
weight.

\paragraph{Contributions.}
\begin{enumerate}[leftmargin=1.4em,itemsep=1pt,topsep=2pt]
\item \textbf{Grid-aware sub-native compression.} To our knowledge, GRAFT is the first method that compresses
FP4-native checkpoints below 4 bits per weight by exploiting their grid. Prior work on these checkpoints stays at about
4 bits and treats the weights as continuous~\citep{grella2026trit,requant2026}.
\item \textbf{A discrete-source rate--distortion analysis} of block-scaled FP4 weights. At 3 bits the memoryless
distortion--rate function of DeepSeek's weights is about \final{DS-BOUND3-1}$\times$ that of a Gaussian source, for
which Gaussian-matched rotation-based coders are designed. The source is close to memoryless in our measurements,
so we use the i.i.d.\ function as a reference (\cref{sec:theory}).
\item \textbf{A trellis code that decodes to the source's own format.} Grid-aligned trellis quantizers exist for
uniform integer grids~\citep{marcellin1990tcq,nnr2022,ccq2025} and for FP4 values~\citep{que2026xortrellis}.
GRAFT decodes to OCP MXFP4 blocks on the grid of the checkpoint it compresses, with power-of-two block scales, a
per-column output scale and a sparse correction. The block-scaled product therefore runs on
existing MX tensor-core GEMMs, including the reference model's own expert kernel, with only the output scale and the
sparse correction added (\cref{sec:method}).
\item \textbf{Results on DeepSeek-V4.1-Flash} for all 15,360 routed experts, compared with EXL3 at 3 bits at the
expert level and at the model level (excess NLL, KL and top-1 agreement on held-out text). Ablations separate the
code from the fitting objective by refitting EXL3 with GRAFT's objective, and they compare block flags with a
per-matrix octave codebook, a flat rate with a per-matrix rate allocation, and the trellis with 3-bit scalar codebooks
(\cref{sec:experiments}).
\end{enumerate}
\ifdraftmarks\noindent\textit{Draft status.} Blue numbers exist and are preliminary; red boxes mark the model-level
evaluation of GRAFT in progress.\fi
